\section{Conclusion}
\label{sec:conclusion}

This paper provides the first direct comparison of the \ar{} algorithmic
redistricting plan with published \recom{} ensemble distributions for six
congressional states.  The results are consistent and interpretable:

\begin{enumerate}
  \item The \ar{} plan is more compact than the average valid plan in
    every state studied, falling at the $61$st--$72$nd percentile of
    Polsby-Popper compactness.  This reflects the METIS edge-cut
    objective, which directly minimises the number of shared census-tract
    boundaries across district lines.

  \item The \ar{} plan produces partisan outcomes near the ensemble median
    in 5 of 6 states (54th, 61st, 52nd, 52nd, and 55th percentiles for
    NC, WI, PA, TX, and CA respectively).  Georgia is the exception at
    the 38th percentile, reflecting the state's strong geographic partisan
    sorting rather than any algorithmic bias.

  \item Minority-majority district counts are consistent with ensemble
    medians in all states with published data, confirming that the \ar{}
    plan does not sacrifice minority representation for compactness.
\end{enumerate}

The legal implication is direct: the \ar{} plan can be presented to courts
with two certificates.  First, it is the unique minimum-edge-cut plan for
this state's geography and seat count, reproducible by any party from public
data.  Second, it lies in the interior of the space of geographically
reasonable plans on all three dimensions that courts have examined.

Future work will extend this comparison to state legislative maps (F-series)
and will compute bespoke \recom{} ensembles for Texas and California to
replace the literature bounds used here.
